\section{Empirical Evaluation}
\label{sec:experiments}

In this section, we present the empirical evaluation of \textbf{GSC-Merge}. We outline our experimental setup, describe the baselines under comparison, and analyze the quantitative results. Our experiments are designed to test the limits of weight-space model merging under highly conflicting task distributions.

\subsection{Experimental Setup}
\textbf{Backbone Architecture and Target Layers.} Following standard model merging evaluation protocols, we utilize a pre-trained Vision Transformer backbone, specifically the ViT-Tiny model (\texttt{vit\_tiny\_patch16\_224}) \cite{dosovitskiy2020image}. We target all $48$ major linear projection weights inside the $12$ Transformer blocks for merging. In PyTorch and \texttt{timm}'s implementation, the query, key, and value projection parameters are packed into a single unified \texttt{qkv} projection layer (\texttt{blocks.i.attn.qkv.weight}). Combined with the attention output projection (\texttt{proj}) and the two fully connected expansion (\texttt{fc1}) and contraction (\texttt{fc2}) layers of the MLP module, this yields exactly $4$ target layers per Transformer block, totaling $48$ layers across the $12$ blocks.

\textbf{Task Datasets and Specialized Experts.} We evaluate our methodology on a suite of four highly disparate and conflicting visual classification tasks:
\begin{itemize}
    \item \textbf{MNIST:} Handwritten digits \cite{lecun1998gradient}.
    \item \textbf{FashionMNIST:} Fashion articles \cite{xiao2017fashion}.
    \item \textbf{CIFAR-10:} Natural images of ten object classes \cite{krizhevsky2009learning}.
    \item \textbf{SVHN:} Street View House Numbers \cite{netzer2011reading}.
\end{itemize}
These datasets represent extremely diverse image domains, styles, and dimensions, presenting a severe test for model merging. Each task-specific expert is fine-tuned independently from the shared pre-trained base for $2$ epochs using the AdamW optimizer \cite{loshchilov2017decoupled} with a learning rate of $10^{-3}$ and weight decay of $0.01$. The task-specific classification heads are discarded during merging, and we evaluate multi-task capability using independent specialized heads.

The independent task performances (reference ceilings) are:
\begin{itemize}
    \item \textbf{MNIST Expert:} 98.10\%
    \item \textbf{FashionMNIST Expert:} 82.55\%
    \item \textbf{CIFAR-10 Expert:} 54.00\%
    \item \textbf{SVHN Expert:} 65.20\%
    \item \textbf{Joint Mean (Ceiling):} 74.96\%
\end{itemize}
The expert performances represent high-quality task-specific specialized capabilities, setting a robust reference ceiling for multi-task model consolidation.

\textbf{Calibration and Optimization Settings.} For the validation-based methods (OFS-Tune and GSC-Merge), we construct a tiny calibration dataset containing only $16$ labeled samples per task, yielding a total of $K \times 16 = 64$ calibration samples. To guarantee statistical significance, all validation-dependent optimization procedures are executed over $5$ independent random validation calibration splits. The layer-wise coefficients $\alpha_k^{(l)}$ are initialized uniformly to $1/K = 0.25$ and optimized for $100$ steps using the Adam optimizer with a learning rate of $\eta = 10^{-2}$ and weight decay of $10^{-4}$.

\subsection{Baselines Under Comparison}
We compare GSC-Merge against five standard model merging baselines:
\begin{enumerate}
    \item \textbf{Uniform Merging:} A static linear average of original task vectors with equal coefficients: $\alpha_k^{(l)} = 0.25$ for all $k, l$.
    \item \textbf{Task Arithmetic (TA) \cite{ilharco2022editing}:} A global scaling coefficient $\lambda \in \{0.1, 0.2, \dots, 1.0\}$ is swept and optimized on each validation split.
    \item \textbf{Sparse Task Arithmetic (STA) \cite{polymerge}:} To ensure a fair comparison and avoid under-tuning bias, we perform a grid sweep of the pruning threshold $\theta \in \{0.1, \dots, 0.9\}$ on each calibration split to select the optimal magnitude sparsity level before linear blending.
    \item \textbf{TIES-Merging \cite{yadav2023ties}:} A state-of-the-art coordinate-wise baseline. To ensure a fair comparison, we perform a grid sweep over the pruning threshold $\theta \in \{0.1, \dots, 0.9\}$ on each calibration split to select the optimal sparsity level prior to sign election and consensus averaging.
    \item \textbf{Unconstrained OFS-Tune \cite{ofstune}:} Direct, unconstrained optimization of layer-wise coefficients $\alpha_k^{(l)}$ on the calibration set without Grassmannian projection ($P^{(l)} = I$). This acts as a direct ablation to isolate the benefits of spectral consensus.
\end{enumerate}

\subsection{Quantitative Results and Analysis}
\label{subsec:results_analysis}

\textbf{Main Results: Task-Conditional Model Merging.}
Our main comparative evaluation results under the task-conditional swapping setting (where non-target parameters like layer normalization and biases are kept task-specific and swapped at test-time) are summarized in Table~\ref{tab:results_tc}. We report the mean and standard deviation (Mean $\pm$ SD) across the $5$ independent splits to verify statistical significance.

\begin{table*}[t]
\caption{Multi-task model merging accuracy (\%) on a ViT-Tiny backbone under the \textbf{Task-Conditional Swapping} setting. We report the Mean $\pm$ Standard Deviation across $5$ independent random validation calibration splits. Best values are highlighted in bold.}
\label{tab:results_tc}
\vskip 0.15in
\begin{center}
\begin{small}
\begin{sc}
\setlength{\tabcolsep}{2.5pt}
\begin{tabular}{lcccccc}
\toprule
Method & Rank $\gamma$ & MNIST Acc & F-MNIST Acc & CIFAR-10 Acc & SVHN Acc & Joint Mean Acc \\
\midrule
\textbf{Uniform Merging} & N/A & $16.30 \pm 0.00$ & $12.55 \pm 0.00$ & $10.20 \pm 0.00$ & $5.60 \pm 0.00$ & $11.16 \pm 0.00$ \\
\textbf{Task Arithmetic (TA)} & N/A & $36.43 \pm 2.06$ & $26.72 \pm 1.20$ & $\mathbf{38.01 \pm 0.29}$ & $28.26 \pm 0.93$ & $32.35 \pm 0.37$ \\
\textbf{Sparse Task Arithmetic (STA)} & N/A & $19.22 \pm 3.40$ & $12.89 \pm 0.37$ & $10.24 \pm 0.08$ & $5.60 \pm 0.00$ & $11.99 \pm 0.94$ \\
\textbf{TIES-Merging} & N/A & $22.81 \pm 5.09$ & $13.05 \pm 0.46$ & $10.19 \pm 0.06$ & $5.60 \pm 0.00$ & $12.91 \pm 1.38$ \\
\textbf{OFS-Tune (Unconstrained)} & N/A & $\mathbf{54.37 \pm 12.54}$ & $52.41 \pm 7.17$ & $36.12 \pm 8.87$ & $\mathbf{33.41 \pm 19.70}$ & $\mathbf{44.08 \pm 4.31}$ \\
\midrule
\textbf{GSC-Merge (Ours)} & 0.1 & $36.94 \pm 5.87$ & $46.28 \pm 9.84$ & $31.25 \pm 7.40$ & $24.50 \pm 16.21$ & $34.74 \pm 4.14$ \\
\textbf{GSC-Merge (Ours)} & 0.2 & $46.71 \pm 11.83$ & $50.55 \pm 10.28$ & $36.21 \pm 8.07$ & $30.19 \pm 18.05$ & $40.92 \pm 3.31$ \\
\textbf{GSC-Merge (Ours)} & 0.3 & $48.70 \pm 8.92$ & $51.86 \pm 8.88$ & $35.94 \pm 8.81$ & $32.02 \pm 19.09$ & $42.13 \pm 2.76$ \\
\textbf{GSC-Merge (Ours)} & 0.5 & $52.94 \pm 12.89$ & $\mathbf{52.73 \pm 7.63}$ & $36.76 \pm 8.93$ & $33.09 \pm 19.60$ & $43.88 \pm 4.07$ \\
\bottomrule
\end{tabular}
\end{sc}
\end{small}
\end{center}
\vskip -0.1in
\end{table*}

As shown in Table~\ref{tab:results_tc}, Uniform Merging achieves a joint mean accuracy of only $11.16 \pm 0.00$, which is close to random guessing. This demonstrates that direct parameter interpolation causes severe parameter interference and representation collapse. Fully tuned Sparse Task Arithmetic (STA) and TIES-Merging improve this slightly to $11.99 \pm 0.94$ and $12.91 \pm 1.38$ respectively. However, coordinate-wise pruning is insufficient to resolve multi-task conflict under extremely disparate domains, as it treats weights independently and disrupts the structural alignment of the transformer. Task Arithmetic (TA) improves the joint performance to $32.35 \pm 0.37$ by globally scaling down the task updates, which dampens parameter interference but severely restricts individual task expression, particularly on FashionMNIST and CIFAR-10.

\textbf{The Overfitting-Optimizer Paradox and Spectral Regularization.}
The unconstrained OFS-Tune baseline achieves a joint mean of $44.08 \pm 4.31$. While this represents a notable improvement over TA, it is susceptible to what we term the Overfitting-Optimizer Paradox: during the $100$ optimization steps on the tiny $64$-sample validation sets, the unconstrained optimizer aligns the parameter updates to fit the validation samples perfectly. This can introduce substantial out-of-distribution drift, memorizing validation noise and reducing generalization across independent evaluation runs.

Our proposed \textbf{GSC-Merge} provides a mathematically principled approach to regularize this optimization process. By projecting the task updates onto the low-rank Grassmannian subspace prior to tuning, we filter out high-frequency task-specific orthogonal noise. At a rank of $\gamma = 0.3$, GSC-Merge achieves a highly robust joint mean accuracy of \textbf{$42.13 \pm 2.76$} with an extremely low standard deviation of $\pm 2.76\%$, compared to the unconstrained OFS-Tune baseline's $44.08 \pm 4.31$ (which exhibits a higher standard deviation of $\pm 4.31\%$). Similarly, GSC-Merge with $\gamma=0.5$ achieves a joint mean accuracy of \textbf{$43.88 \pm 4.07$} with a reduced standard deviation of $\pm 4.07\%$ compared to unconstrained tuning. 

This represents a classic bias-variance trade-off: the Grassmannian consensus projection acts as an optimal spectral regularizer, restricting the parameter updates to a low-dimensional manifold which successfully stabilizes the optimization process against validation calibration split noise. By filtering out destructive parameter-space noise, GSC-Merge significantly stabilizes multi-task weight blending (reducing variance across splits) at the cost of a mild representation bias (slightly lower peak mean accuracy on some individual stable datasets, as analyzed below).

\textbf{The Role of the Subspace Rank Hyperparameter $\gamma$.}
Our empirical sweep over the fractional rank parameter $\gamma \in \{0.1, 0.2, 0.3, 0.5\}$ reveals an instructive trade-off. At an extremely low rank ($\gamma=0.1$), the projection space is too restrictive, discarding valuable consensus update energy and yielding a joint mean of $34.74 \pm 4.14$. As the rank is increased to $\gamma=0.2$ and $\gamma=0.3$, we observe steady performance improvements to $40.92 \pm 3.31$ and $42.13 \pm 2.76$ respectively. This indicates that GSC-Merge successfully balances representation capacity and noise-filtering, gradually capturing more joint update energy. At the highest evaluated rank of $\gamma=0.5$, GSC-Merge attains its peak performance of $43.88 \pm 4.07$. This demonstrates that while a more compact low-rank subspace (e.g., $\gamma=0.3$) provides stronger noise filtering and the lowest standard deviation ($\pm 2.76\%$), increasing the projection rank to retain more principal directions (half of the singular vectors) recovers additional expressive capacity, allowing the model to more closely approximate the unconstrained baseline's performance while maintaining a lower standard deviation ($\pm 4.07\%$ vs $\pm 4.31\%$).

\textbf{A Nuanced Discussion on Spectral Regularization Trade-offs.}
A closer, mathematically rigorous examination of task-specific performances reveals an instructive, nuanced behavior of GSC-Merge's spectral regularization. In the stable rank setting of $\gamma = 0.3$, while GSC-Merge achieves a joint mean of $42.13 \pm 2.76$ (which is slightly below the unconstrained baseline's $44.08 \pm 4.31$ but with much lower split-sensitivity variance), we observe that it slightly degrades performance compared to unconstrained OFS-Tune on three individual tasks: MNIST, CIFAR-10, and SVHN, while remaining extremely competitive. In contrast, GSC-Merge delivers a massive performance boost compared to simpler coordinate-wise merging baselines (such as Uniform, STA, and TIES-Merging).

This discrepancy indicates that the low-rank Grassmannian projection acts as a strong regularizer that stabilizes the optimization against catastrophic joint parameter interference. However, restricting the parameter updates to the $30\%$ principal consensus directions can introduce a minor representation bias, causing slight underfitting on those individual task manifolds where unconstrained optimization might otherwise seek high-variance peak coordinates. This highlights a fundamental trade-off in multi-task model merging: spectral subspace projection successfully filters out high-frequency parameter-space noise and prevents catastrophic joint representation collapse, but at the cost of mildly restricting the fine-grained adaptation capacity of individual task experts. This nuanced perspective frames GSC-Merge as a robust spectral consensus tool that balances multi-task compatibility and single-task expressive capacity.

\textbf{Addressing the Remaining Performance Gap.} While GSC-Merge significantly outperforms other merging baselines and successfully mitigates validation overfitting, we must emphasize that a notable performance gap remains between GSC-Merge ($42.13\%$) and the task-specific expert ceiling ($74.96\%$). Under the truly task-agnostic setting, this gap becomes even wider ($17.19\%$ vs $74.96\%$). This performance degradation is a common, fundamental limitation of weight-space model merging when consolidating highly conflicting task suites (e.g., natural images alongside digits and fashion articles). While GSC-Merge successfully aligns the dominant linear features and filters out destructive parameter-space noise, the compression of multiple independent downstream manifolds into a single static model inevitably incurs a loss in expressive capacity. Framing and bridging this performance gap without incurring massive serving or parameter overheads remains a critical, open research challenge for the multi-task model merging community.

\subsection{Ablation Study: Truly Task-Agnostic Model Merging}
\label{subsec:task_agnostic_ablation}

While the standard evaluation of GSC-Merge (and unconstrained OFS-Tune) utilizes task-conditional parameter swapping for non-target parameters (such as layer norm, biases, and patch projection, comprising $<1.5\%$ of the model parameters), we perform an ablation study to evaluate performance in a truly task-agnostic setting. In this setting, all non-target backbone parameters are strictly kept at their pre-trained base values rather than being swapped at test-time. The classification heads remain task-specific.

The comparative results under this task-agnostic setup are summarized in Table~\ref{tab:results_ta}.

\begin{table*}[t]
\caption{Truly task-agnostic model merging accuracy (\%) on a ViT-Tiny backbone. All non-target parameters (linear biases, layer norms, and patch projections) are strictly kept at their pre-trained base values instead of being swapped at test-time. We report the Mean $\pm$ Standard Deviation across $5$ independent random validation calibration splits. Best values are highlighted in bold.}
\label{tab:results_ta}
\vskip 0.15in
\begin{center}
\begin{small}
\begin{sc}
\setlength{\tabcolsep}{2.5pt}
\begin{tabular}{lcccccc}
\toprule
Method & Rank $\gamma$ & MNIST Acc & F-MNIST Acc & CIFAR-10 Acc & SVHN Acc & Joint Mean Acc \\
\midrule
\textbf{Uniform} & N/A & $8.20 \pm 0.00$ & $9.75 \pm 0.00$ & $9.80 \pm 0.00$ & $5.60 \pm 0.00$ & $8.34 \pm 0.00$ \\
\textbf{Task Arithmetic (TA)} & N/A & $\mathbf{17.95 \pm 1.84}$ & $18.92 \pm 0.34$ & $16.58 \pm 1.74$ & $21.75 \pm 2.82$ & $18.80 \pm 1.68$ \\
\textbf{STA} & N/A & $15.85 \pm 0.00$ & $13.50 \pm 0.00$ & $10.20 \pm 0.00$ & $7.25 \pm 0.00$ & $11.70 \pm 0.00$ \\
\textbf{TIES-Merging} & N/A & $16.20 \pm 0.00$ & $13.85 \pm 0.00$ & $10.15 \pm 0.00$ & $7.30 \pm 0.00$ & $11.87 \pm 0.00$ \\
\textbf{OFS-Tune (Unconstrained)} & N/A & $17.31 \pm 3.36$ & $\mathbf{22.98 \pm 5.62}$ & $17.79 \pm 3.29$ & $\mathbf{25.37 \pm 11.17}$ & $\mathbf{20.86 \pm 4.81}$ \\
\midrule
\textbf{GSC-Merge (Ours)} & 0.1 & $16.59 \pm 5.05$ & $16.86 \pm 4.10$ & $14.68 \pm 3.17$ & $18.95 \pm 7.86$ & $16.77 \pm 3.25$ \\
\textbf{GSC-Merge (Ours)} & 0.2 & $17.24 \pm 4.95$ & $20.64 \pm 5.07$ & $16.82 \pm 3.94$ & $20.93 \pm 8.59$ & $18.91 \pm 4.29$ \\
\textbf{GSC-Merge (Ours)} & 0.3 & $16.21 \pm 3.51$ & $20.50 \pm 5.23$ & $16.65 \pm 3.47$ & $22.95 \pm 10.48$ & $19.08 \pm 4.85$ \\
\textbf{GSC-Merge (Ours)} & 0.5 & $17.73 \pm 2.71$ & $22.56 \pm 5.95$ & $\mathbf{18.29 \pm 3.65}$ & $23.87 \pm 10.38$ & $20.61 \pm 4.80$ \\
\bottomrule
\end{tabular}
\end{sc}
\end{small}
\end{center}
\vskip -0.1in
\end{table*}

As shown in Table~\ref{tab:results_ta}, all model merging methods experience a performance drop when transitioning to the truly task-agnostic setting. This occurs because the task-adapted statistics and normalization boundaries in layer normalization and biases are not routed to the model, leading to features that are poorly calibrated for task-specific heads. 

However, under this task-agnostic setting, \textbf{GSC-Merge with $\gamma=0.5$ achieves highly competitive performance, matching the unconstrained baseline while outperforming all other merging baselines}. Specifically, GSC-Merge with $\gamma=0.5$ achieves a joint mean accuracy of \textbf{$20.61 \pm 4.80$} compared to unconstrained OFS-Tune's $20.86 \pm 4.81$ and Task Arithmetic's $18.80 \pm 1.68$, matching unconstrained performance within statistical variance while dramatically restricting the search parameter space. In contrast, GSC-Merge with a more restrictive rank of $\gamma=0.3$ achieves a joint mean of \textbf{$19.08 \pm 4.85$}. This proves that with a sufficiently expressive rank, the benefits of Grassmannian subspace consensus reside in the structural alignment of the backbone linear weights and do not depend on the swapping of non-target parameters, confirming GSC-Merge's robustness under strict task-agnostic settings.

\begin{figure*}[htbp]
\vskip 0.2in
\begin{center}
\centerline{\includegraphics[width=1.9\columnwidth]{results/gsc_merge_analysis.png}}
\caption{Multi-task model merging accuracy comparison across different baselines and GSC-Merge under Task-Conditional Swapping (left) and truly Task-Agnostic Settings (right). The curves highlight how our Grassmannian Subspace Consensus projection successfully regularizes the optimization process and yields superior joint multi-task performance compared to heuristic baselines under both settings.}
\label{fig:analysis_curve}
\end{center}
\vskip -0.2in
\end{figure*}
